\documentclass[11pt,a4paper]{article}
\usepackage{float}

\usepackage[utf8]{inputenc}
\usepackage[T1]{fontenc}
\usepackage[english]{babel}

\usepackage{amsmath,amssymb,amsthm}
\usepackage{mathtools}
\usepackage{mathrsfs}
\usepackage{graphicx}
\usepackage{hyperref}
\usepackage{geometry}
\usepackage{physics}
\usepackage{enumitem}
\setlist[itemize,1]{label=---}
\geometry{margin=2.5cm}
\usepackage{csquotes}

\newtheorem{theorem}{Theorem}[section]
\newtheorem{proposition}[theorem]{Proposition}
\newtheorem{lemma}[theorem]{Lemma}
\newtheorem{corollary}[theorem]{Corollary}

\theoremstyle{definition}
\newtheorem{definition}[theorem]{Definition}
\newtheorem{remark}[theorem]{Remark}
\newtheorem{example}[theorem]{Example}

\newcommand{\PDL}{\mathrm{PDL}}
\newcommand{\LDP}{\mathrm{LDP}}

\begin{document}

\title{Whoever We May Be:\\[4pt]
Existence, Coherence, and Vulnerable Instruments\\[4pt]
A Philosophical Synthesis of the Projective Dynamic Logo Framework}

\author{C.~Laubscher\thanks{Independent researcher, Switzerland.}\\[4pt]}

\date{\small \today}

\maketitle

\begin{abstract}
This paper offers a philosophical and ontological synthesis of the Projective Dynamic Logo (PDL) framework, intended as an accessible point of entry for both practicing scientists and intellectually engaged non-specialists. It takes as its guiding thread the claim that whatever we may be---from subatomic closures to conscious agents---we are the vulnerable instruments of a simple yet relentless logic of coherence, formalised in PDL by axioms on finite signed graphs and the selection of pulsating closures. Rather than introducing new technical results, the text reinterprets the existing PDL corpus as a proposal about what it means to exist in a logically constrained universe, clarifying how notions of causality, freedom, transcendence, and human fragility can be read, with appropriate caution, against the background of minimal closure, coherence leakage, and hierarchical architectures already developed in the technical literature.
\end{abstract}

\tableofcontents

\section{Introduction: A Sentence and a Programme}
\label{sec:intro}

\subsection{Whoever we may be}

This paper is guided by a single sentence, which first emerged not as a slogan, but as a condensed expression of a long-standing intuition:

\begin{quote}
Whoever we may be, whatever we may be, from the atoms that form our molecules and cells to our very thoughts and souls, we are but the necessary expression of a simple and relentless logic of coherence, which leaves no event beyond its reach. A logic of causality that fulfils itself with every pulse. It is the Origin, the engine, and the essence of the universe. The human being is nothing more than its most vulnerable instrument.
\end{quote}

Read in isolation, this statement may appear merely poetic or metaphysical. Within the Projective Dynamic Logo (PDL) framework, however, it can be understood as an attempt to translate into ordinary language what a specific family of axioms and constructions already articulate in a rigorous, if austere, manner.\cite{Laubscher_IntroPDL_2024,Laubscher_EmergencePhysicalReality_2024,Laubscher_PDL_Position_2024} PDL does not begin with space, time, or matter; it begins with reproducible distinctions on finite signed graphs, with minimal pulsating closures that are selected by coherence-oriented constraints rather than imposed by hand.

The central claim of the present text is not that this sentence is true in some final metaphysical sense, but that it faithfully captures the ontological stance implicit in the technical PDL corpus, once that corpus is read as more than a bare mathematical exercise. The task here is therefore exegetical and interpretative: to show how a programme of discrete graph-theoretic constructions about electrons, protons, constants, and emergent metrics can be viewed as a proposal about what it \emph{means} to exist in a logically constrained universe.\cite{Laubscher_Existence_PulsatingClosure_2026,Laubscher_PDL_GlobalMapping_v6_2024}

\subsection{From technical austerity to ontological stakes}

Most existing PDL papers are deliberately written in a technically minimalist and conceptually sparse style.\cite{Laubscher_EmergencePhysicalReality_2024,Laubscher_MinimalClosures_46_2024,Laubscher_ProtonArch_v2_2024,Laubscher_Alpha_v2_2024,Laubscher_EffectiveFields_2024,Laubscher_CavityModes_PDL_2024} They define finite signed graphs, state axioms of binary pulsation and triangular coherence, derive the necessity of the $(4,6)$ block as the first admissible stationary closure, construct a rigid proton architecture specified by a small list of integers, and propose combinatorial expressions for quantities such as the fine-structure constant, Boltzmann's constant, and an effective gravitational coupling. The tone is intentionally dry: proofs, constructions, and numerical comparisons take precedence over narrative or interpretation.

At the same time, scattered across this corpus are conceptual gestures that point beyond pure technicality. Position papers and global mappings explicitly describe PDL as a \emph{foundational research programme} concerned with the conditions under which persistent existence can be ascribed at all.\cite{Laubscher_PDL_Position_2024,Laubscher_PDL_GlobalMapping_v6_2024} The ontological meditation on ``existence as pulsating closure'' makes clear that the move to finite signed graphs is not a mere computational convenience, but a proposal about the nature of reality: to exist is to instantiate a self-sustaining cycle of coherence on a bounded relational structure.\cite{Laubscher_Existence_PulsatingClosure_2026} In other words, the technical austerity of the PDL texts conceals, without erasing, strong ontological commitments.

The aim of this paper is to make those commitments explicit in a manner that is accessible both to practising scientists and to non-specialists who are prepared to engage with conceptual arguments, without requiring them to follow every combinatorial proof. To that end, the text is written at two levels. The main body develops the central ideas in ordinary language, with only light use of notation. A parallel layer of ``technical anchors''---short, clearly marked passages and references---indicates where each conceptual claim is rooted in specific PDL constructions and theorems. Readers primarily interested in the vision may skim the anchors; readers concerned with rigour may treat them as pointers back into the formal corpus.

\subsection{Scope, audience, and epistemic status}

It is important to delineate, from the outset, the scope and epistemic status of what follows. This paper does not introduce new theorems, new numerical derivations, or new empirical predictions within the PDL framework. All technical results cited here have been presented elsewhere: the emergence of the $(4,6)$ block as minimal stationary closure, the combinatorial selection of the proton architecture, the reinterpretation of the fine-structure constant and related quantities, the analysis of coherence leakage and the exponent $18$, the construction of effective fields and cavity modes, and the proposal of an Einstein--Dirac interface.\cite{Laubscher_MinimalClosures_46_2024,Laubscher_ProtonArch_v2_2024,Laubscher_Alpha_v2_2024,Laubscher_Exponent18_v2_2024,Laubscher_EffectiveFields_2024,Laubscher_CavityModes_PDL_2024,Laubscher_Einstein_Dirac_2024}

Nor does the paper claim to offer a complete theory of life, consciousness, or value. Any attempt to deduce the richness of human experience from finite signed graphs alone would not only outstrip the current state of PDL, but also risk an implausible reductionism. Where notions such as freedom, vulnerability, or transcendence are discussed, they are presented as \emph{interpretative extrapolations}: ways of reading familiar phenomena in the light of an ontology of pulsating closures and incomplete structures, not as derivations or proofs.\cite{Laubscher_Existence_PulsatingClosure_2026,Laubscher_PDL_TO_Dialogue_2024}

The target audience is deliberately mixed. On the one hand, the paper addresses researchers in the philosophy and foundations of physics who wish to understand what PDL is really saying about existence, beyond its technical apparatus. On the other hand, it seeks to remain readable to scientifically literate non-specialists who are willing to encounter some unfamiliar terminology in exchange for a coherent picture. Navigating this dual audience requires a certain discipline. Whenever a concept is introduced that has a precise formal counterpart in the PDL corpus, the relation between the informal and formal versions will be stated as clearly as possible, and references will be provided. Whenever the discussion moves into speculative or existential territory, this will be explicitly signalled.

Within this framework, the guiding sentence quoted above can be reformulated more cautiously as a working hypothesis: \emph{if} PDL offers a viable account of how physical reality is organised, \emph{then} it becomes meaningful to regard ourselves---and every other persistent structure---as finite, vulnerable instruments through which a logically constrained coherence is realised and partially reflected. The remainder of the paper is devoted to unpacking what this hypothesis entails.

\section*{Glossary of Key Terms}

\begin{description}

\item[Finite signed graph]
A finite set of vertices connected by edges, each edge carrying a sign $+1$ or $-1$. In PDL this is the basic substrate on which all closures, particles, and fields are defined.\cite{Laubscher_IntroPDL_2024,Laubscher_EmergencePhysicalReality_2024}

\item[Closure]
A relational pattern on a finite signed graph that approximately satisfies prescribed coherence conditions on its triangular subgraphs and can sustain a recurrent internal dynamics. In PDL, to ``exist'' in the minimal sense is to realise such a closure.\cite{Laubscher_EmergencePhysicalReality_2024,Laubscher_MinimalClosures_46_2024}

\item[Minimal stationary closure]
A closure that (i) is defined on a graph that cannot be reduced further without losing coherence, and (ii) supports a stationary pulsating regime under the update rule. The first such closure in PDL is the $(4,6)$ block.\cite{Laubscher_MinimalClosures_46_2024}

\item[$(4,6)$ block]
The complete signed graph on four vertices and six edges, shown to be the first admissible minimal stationary closure under the PDL axioms. It is interpreted as the logical prototype of an electron at rest.\cite{Laubscher_MinimalClosures_46_2024,Laubscher_EmergencePhysicalReality_2024}

\item[Pulsating closure]
A closure whose internal sign configuration evolves under a discrete update rule and returns to itself after a finite number of steps (often a logical $2$-cycle). The repeated restoration of coherence in this cycle is the formal counterpart of a ``pulse'' in the guiding sentence.\cite{Laubscher_IntroPDL_2024,Laubscher_Existence_PulsatingClosure_2026}

\item[Triangular coherence]
A local closure condition imposed on each triangle (three mutually connected vertices) of a signed graph, typically formulated as a constraint on the product of its edge signs. The global coherence of a configuration is assessed by counting coherent versus violated triangles.\cite{Laubscher_EmergencePhysicalReality_2024,Laubscher_Leakage_Probability_2024}

\item[Coherence leakage ($\varepsilon$)]
The fraction of violated triangles in an optimised configuration, i.e.\ the ratio of local closure requirements that cannot be satisfied simultaneously. For composite structures $\varepsilon$ is small but non-zero and plays a central role in probabilistic behaviour and emergent couplings.\cite{Laubscher_Leakage_Probability_2024,Laubscher_Exponent18_v2_2024}

\item[Coherence budget]
An informal term for the total amount of coherence a structure can maintain across its internal relations, its active surfaces, and its leakage into the environment. Mass, energy, and certain constants can be interpreted as coarse summaries of how this budget is allocated.\cite{Laubscher_EmergencePhysicalReality_2024,Laubscher_Alpha_v2_2024,Laubscher_Existence_PulsatingClosure_2026}

\item[Active surface]
The subset of a composite architecture’s relations that effectively couple it to external closures. In the PDL proton model, the active surface is the part of the relational budget through which an electron-like $(4,6)$ block can interact with the proton.\cite{Laubscher_ProtonArch_v2_2024,Laubscher_Born_Golden_v2_2024}

\item[Proton architecture]
A hierarchically structured composite built from $(4,6)$ blocks under additional axioms on valence structure, connectivity, and coherence. It is characterised by specific integers (such as $24,28,930,10087,11017$) and underpins structural proposals for physical constants.\cite{Laubscher_ProtonArch_v2_2024,Laubscher_PDL_GlobalMapping_v6_2024}

\item[Hierarchical filtering]
The process by which local coherence leakage at small scales (e.g.\ nucleonic) is propagated through a sequence of higher-level structures, amplifying a tiny $\varepsilon$ into an effective macroscopic coupling. In PDL this is modelled with a minimal number of filters, often associated with an exponent around $18$.\cite{Laubscher_Exponent18_v2_2024,Laubscher_UniversalCoherenceLeakage_2024}

\item[Effective field]
A coarse-grained description of how coherence deficits and surpluses are redistributed across many closures, appearing at a large scale as a field-like degree of freedom (for example, an emergent metric or force field).\cite{Laubscher_EffectiveFields_2024}

\item[Mode and density of states]
In PDL cavity models, a mode is a periodic orbit of a photon-like excitation on a finite signed graph, indexed by a triple of integers. The density of states describes how the number of such modes grows with logical frequency, reproducing the familiar $\rho(\nu)\propto \nu^2$ scaling in three dimensions.\cite{Laubscher_CavityModes_PDL_2024}

\item[Transcendence (in PDL)]
Not a separate realm beyond physics, but the structural fact that no finite closure can fully absorb its own coherence requirements. Incomplete closure and non-zero leakage mean that every entity depends on and refers beyond itself.\cite{Laubscher_Existence_PulsatingClosure_2026,Laubscher_PDL_TO_Dialogue_2024}

\item[Vulnerable instrument]
A phrase used to describe human beings as high-level closures that are both capable of reflecting on their own coherence management and exposed to leakage and dependence on wider relational structures. Vulnerability here is the lived manifestation of incomplete closure.\cite{Laubscher_Existence_PulsatingClosure_2026}

\end{description}

\section{Minimal Coherence: From Nothingness to Pulsation}
\label{sec:minimal_coherence}

\subsection{Suspending space, time, and matter}

The PDL framework begins with a methodological act of suspension. Instead of taking for granted a continuous spacetime manifold populated by particles and fields, it brackets these familiar ingredients and asks a more austere question: \emph{what is the least amount of structure required for there to be something rather than nothing, in a way that can persist}? \cite{Laubscher_IntroPDL_2024,Laubscher_EmergencePhysicalReality_2024,Laubscher_Existence_PulsatingClosure_2026} In this upstream setting, one cannot appeal to distances, durations, or energies, because the very notions of metric, trajectory, and dynamical law are precisely what is under scrutiny.

From this perspective, neither an instantaneous flash nor a perfectly featureless state qualifies as existence in a robust sense. An event that occurs only once, without any possibility of recurrence or recognition, is empirically and conceptually indistinguishable from its non-occurrence; a homogeneous, unchanging background offers no internal contrasts by which anything could be said to be present rather than absent.\cite{Laubscher_Existence_PulsatingClosure_2026} Existence, in the sense at stake here, must therefore involve a \emph{reproducible distinction}: some pattern or difference that can return to itself and be re-instantiated.

PDL formalises this intuition by replacing the continuum with finite signed graphs. Vertices stand for abstract informational loci, edges carry binary signs encoding an elementary attitude of agreement or disagreement, and small subgraphs (triangles, in particular) serve as the minimal arenas in which coherence or incoherence can be evaluated.\cite{Laubscher_EmergencePhysicalReality_2024,Laubscher_PDL_Position_2024} In doing so, the framework shifts the focus from objects in space to relational configurations that must be capable of sustaining their own pattern without leaning on an undefined background.

\subsection{Reproducible distinction and binary pulsation}

At the most elementary level, a reproducible distinction is modelled in PDL as a binary pulsation: an alternation between two mutually exclusive states that can return to itself.\cite{Laubscher_IntroPDL_2024,Laubscher_Existence_PulsatingClosure_2026} Concretely, one considers sign assignments on edges of a finite graph and a discrete update rule that flips or redistributes these signs. If there exists a configuration such that the rule generates a logical \(2\)-cycle---a sequence of two distinct sign patterns that repeat indefinitely---then an elementary pulsating regime has been realised.

At this stage, no physical time parameter is introduced. The ``pulse'' is not yet a vibration in spacetime, but a minimal logical recurrence: a configuration that differs from nothingness by the fact that it can return to itself in a structured way. The idea that ``a logic of causality fulfils itself with every pulse'' is thus given a precise but abstract meaning: causality is not initially a chain of events ordered in time, but the stability of a rule that, when applied repeatedly, maintains a pattern of distinctions instead of letting them decay into undifferentiated noise.\cite{Laubscher_Existence_PulsatingClosure_2026}

This binary pulsation is not presented as a metaphysical axiom about the cosmos; it is a modelling choice about how to represent the minimal conditions for persistent difference. The claim is that any more elaborate conception of existence---from electrons to galaxies---must, at bottom, presuppose the possibility of such recurring distinctions. In the PDL programme, this claim is made testable by deriving concrete structures and constants from it.\cite{Laubscher_EmergencePhysicalReality_2024,Laubscher_PDL_GlobalMapping_v6_2024}

\subsection{Triangles, closure, and logical coherence}

A single pulsating edge, however, is not enough to support a robust notion of existence. An isolated binary alternation, unmoored from any relational context, cannot sustain reference: there is nothing to distinguish one cycle from another, no internal structure that would allow the pattern to be recognised as the ``same'' closure rather than a mere repetition of a trivial flip.\cite{Laubscher_Existence_PulsatingClosure_2026} To avoid smuggling in an implicit background, PDL requires that reproducible distinctions be embedded in small relational circuits that close upon themselves.

In the theory of signed graphs, triangles provide the simplest such circuits. A triangle is the shortest cycle capable of closing without reference to any external vertex, and it already supports a non-trivial notion of balance: depending on the product of its three edge signs, it can be classified as coherent or incoherent.\cite{Laubscher_EmergencePhysicalReality_2024,Laubscher_Leakage_Probability_2024} PDL elevates this observation to an axiom: triangles are the elementary closure motifs, and a global configuration is evaluated by counting how many of these motifs satisfy a prescribed coherence condition.

A configuration is then said to be \emph{logically closed}, in the minimal sense, if it achieves a high degree of triangular coherence without appealing to any structure outside the finite graph on which it is defined.\cite{Laubscher_EmergencePhysicalReality_2024,Laubscher_MinimalClosures_46_2024} Such a configuration does not merely exist as a list of signs; it carries its own criteria of consistency, encoded in the way its subgraphs close. In this way, the intuitive idea of a ``logic of coherence which leaves no event beyond its reach'' is translated into a combinatorial requirement: events that cannot be integrated into closed triangular patterns at an acceptable coherence cost are excluded from minimal existence.

\subsection{Technical anchor: axioms and minimal closures}

Formally, the PDL framework specifies a class of finite signed graphs and a set of qualitative constraints that any admissible minimal closure must satisfy.\cite{Laubscher_IntroPDL_2024,Laubscher_EmergencePhysicalReality_2024,Laubscher_MinimalClosures_46_2024,Laubscher_PDL_Position_2024} The key ingredients can be summarised as follows:

\begin{itemize}
  \item \textbf{Binary pulsation.} There exists a discrete update rule on edge signs such that the induced dynamics admits a non-trivial logical \(2\)-cycle.
  \item \textbf{Triangular coherence.} Triangles are classified as coherent or violated according to a fixed sign product condition; admissible closures must minimise the fraction of violated triangles.
  \item \textbf{Minimal completeness.} The structure must be self-sufficient: no hidden background or larger graph may be presupposed to restore coherence. All necessary relations are contained within the finite signed graph itself.
  \item \textbf{Logical optimisation.} Among graphs that satisfy the preceding qualitative constraints, one favours those that maximise coherence and connectivity while remaining combinatorially minimal.
\end{itemize}

A detailed combinatorial analysis, carried out in dedicated technical work, shows that no finite signed graph with fewer than four vertices can satisfy these constraints simultaneously.\cite{Laubscher_MinimalClosures_46_2024} With one or two vertices, no triangles are present; with three vertices, a single triangle can be rendered coherent, but it cannot sustain a network of overlapping cycles capable of propagating reference without invoking an external hierarchy. The first admissible configuration is the complete signed graph on four vertices and six edges, denoted \((4,6)\), for which one can choose edge signs such that all triangles are coherent and the global sign inversion realises a logical \(2\)-cycle.\cite{Laubscher_MinimalClosures_46_2024,Laubscher_EmergencePhysicalReality_2024}

In PDL, this result is taken as more than a curiosity of graph theory. It is read as an answer to the upstream question posed at the beginning of this section: once spacetime and matter have been suspended, the minimal way for there to be a reproducible something rather than nothing is to instantiate a pulsating closure on a \((4,6)\) signed graph. The next sections explore how this abstract statement acquires physical content when the \((4,6)\) closure is interpreted as the logical prototype of an electron at rest, and how increasingly complex closures give rise to the architectures and constants of familiar physics.\cite{Laubscher_EmergencePhysicalReality_2024,Laubscher_Alpha_v2_2024,Laubscher_ProtonArch_v2_2024}

\section{The First Instrument: The (4,6) Closure and the Electron}
\label{sec:46_electron}

\subsection{From abstract closure to physical prototype}

The \((4,6)\) closure is, at this stage, a purely combinatorial object. It is the first finite signed graph that can sustain a non-trivial logical \(2\)-cycle while satisfying the PDL constraints of triangular coherence, minimal completeness, and optimisation.\cite{Laubscher_MinimalClosures_46_2024,Laubscher_EmergencePhysicalReality_2024} Nothing in its definition mentions charge, mass, spin, or spacetime. Yet, if PDL is to be more than a mathematical curiosity, such closures must eventually be confronted with the phenomenology of actual particles.

Among all known entities in microphysics, the electron at rest presents itself as a natural candidate for this confrontation. It is strictly stable in isolation, exhibits no internal structure at currently accessible length scales, and carries an intrinsic Compton-frequency oscillation associated with its rest mass.\cite{Laubscher_EmergencePhysicalReality_2024,Laubscher_Existence_PulsatingClosure_2026} These features resonate with the logical character of the \((4,6)\) block: a minimal, self-contained pulsating regime that neither decomposes into smaller closures nor relies on a larger background architecture to maintain itself.

The core interpretative step of PDL is therefore a \emph{correspondence principle}: the proposal that the electron at rest should be read as the first physical instantiation of the logically compelled \((4,6)\) closure.\cite{Laubscher_EmergencePhysicalReality_2024,Laubscher_PDL_Position_2024} This is neither a metaphysical identity claim (the electron \emph{is nothing but} a signed graph) nor a complete derivation of quantum electrodynamics. It is a structurally constrained hypothesis: if existence, at its most basic level, consists in sustaining pulsating closures on finite signed graphs, then the simplest such closure ought to be realised somewhere in the empirical world as the simplest available stable microscopic regime.

Once this correspondence is entertained, the guiding sentence of the introduction acquires a first concrete foothold. To say that ``from the atoms that form our molecules and cells to our very thoughts and souls, we are the necessary expression of a logic of coherence'' is to suggest that the logical conditions that pick out the \((4,6)\) closure at the bottom of the PDL hierarchy are not merely abstract constraints, but the first articulation of a structural necessity that percolates through all higher levels of organisation. The electron is then not only a building block of matter, but also the first ``instrument'' through which the underlying logic begins to project itself into a physical regime that we can observe and measure.\cite{Laubscher_Existence_PulsatingClosure_2026,Laubscher_PDL_GlobalMapping_v6_2024}

\subsection{Pulsation, Compton frequency, and the cost of existence}

Once the \((4,6)\) closure is taken as the logical prototype of the electron, its internal pulsation can be compared with the empirical oscillation encoded in the electron’s Compton frequency.\cite{Laubscher_EmergencePhysicalReality_2024,Laubscher_Schrodinger_Compat_v2_2024} In PDL, the logical dynamics on edge signs has period two: at each step of the update rule, the pattern of agreement and disagreement is redistributed, yet after two steps the configuration returns exactly to itself.\cite{Laubscher_MinimalClosures_46_2024} From a purely combinatorial standpoint, this is nothing more than a discrete cycle in a finite state space.

In standard physics, by contrast, the Compton frequency is introduced as a physical parameter: the angular frequency \(\omega_C\) associated with the electron’s rest mass via the relation \(E = \hbar \omega_C\). In conventional treatments this relation is often taken as a compact way of expressing the rest energy of the particle, rather than as a clue to any internal structure. The PDL perspective invites a different reading. If the electron at rest realises the \((4,6)\) closure, then its Compton frequency can be interpreted as the physical manifestation of the internal logical pulsation, now expressed in units of energy and time rather than in purely combinatorial steps.\cite{Laubscher_EmergencePhysicalReality_2024,Laubscher_Existence_PulsatingClosure_2026}

On this view, the reduced Planck constant \(\hbar\) no longer appears as an opaque quantum of action inserted into the formalism from outside. It becomes a measure of the \emph{coherence cost} needed to sustain one full internal cycle of the minimal closure when realised in our universe.\cite{Laubscher_EmergencePhysicalReality_2024,Laubscher_PDL_Position_2024} The familiar equation \(E = \hbar \omega\) is reinterpreted as an accounting identity: it states how much ``coherence budget'' must be expended per unit time to maintain the pulsating regime associated with a given frequency. In this sense, the logical statement that there exists a two-step cycle on the \((4,6)\) graph and the physical statement that the electron has a Compton frequency are seen as two descriptions of the same underlying regularity, written in different languages.

From this vantage point, the guiding idea that there is a ``logic of causality that fulfils itself with every pulse'' ceases to be a mere metaphor. Each cycle of the \((4,6)\) closure can be regarded as a minimal act by which coherence is re-established against the possibility of breakdown, and the associated energy can be read as the price of that act, expressed in physical units.\cite{Laubscher_Existence_PulsatingClosure_2026} To exist, even in the simplest possible way, is already to pay a non-zero coherence cost.

\subsection{Mass, energy, and coherence budgets}

If the reduced Planck constant can be read as the cost of one internal pulsation, it is natural, within PDL, to reinterpret mass and energy as bookkeeping devices for coherence budgets.\cite{Laubscher_EmergencePhysicalReality_2024,Laubscher_Existence_PulsatingClosure_2026} In the conventional picture, mass is a measure of inertia and energy quantifies the capacity to perform work. In the PDL picture, these familiar roles are preserved, but they are rooted in a more primitive notion: the density and organisation of internal coherence cycles that a given structure must sustain in order to persist.

For a minimal closure such as the \((4,6)\) block, the situation is deceptively simple: the internal dynamics is uniform, and the coherence cost per cycle can be associated directly with the rest energy of the corresponding electron prototype via \(E = \hbar \omega_C\).\cite{Laubscher_EmergencePhysicalReality_2024,Laubscher_Schrodinger_Compat_v2_2024} For more complex architectures, such as the proton and higher-level composites built from many \((4,6)\) units, the coherence budget becomes layered. Different parts of the structure---valence cores, relational seas, active surfaces---contribute in distinct ways to the overall cost of maintaining closure, and the resulting mass reflects the combined effort of keeping this multi-scale pattern coherent.\cite{Laubscher_ProtonArch_v2_2024,Laubscher_Alpha_v2_2024}

From this standpoint, inertia can be understood as resistance to changes that would significantly reorganise the internal coherence pattern. A massive object is one whose internal closure is so densely structured that altering its state demands substantial additional coherence investment; a light object corresponds to a relatively sparse or fragile pattern that can be reconfigured at lower cost.\cite{Laubscher_Existence_PulsatingClosure_2026,Laubscher_PDL_Position_2024} This is not intended as a replacement for dynamical equations, but as a reinterpretation of what those equations are \emph{about}. Instead of treating mass and energy as primitive quantities attached to opaque point particles, PDL invites us to see them as summaries of how much and how tightly coherence is being circulated inside underlying relational closures.

In this way, the statement that ``matter appears as a pattern of logical closure so coherent that it has acquired stability and inertia'' becomes more than a suggestive phrase.\cite{Laubscher_Existence_PulsatingClosure_2026} It is a compact description of a concrete proposal: that the traditional parameters of mass and energy encode, in a coarse-grained form, the detailed coherence bookkeeping that PDL attempts to reconstruct explicitly from finite signed graphs.

From this angle, the familiar feeling that some changes are “hard” and others “easy” can be read as a distant echo of this coherence bookkeeping. Rearranging a rigid, tightly organised structure—whether a physical body or a long-held pattern of behaviour—demands more work than adjusting a loose, weakly constrained one. We experience this as inertia, effort, or resistance, but in the PDL picture these are the macroscopic shadows of a much deeper fact: that complex closures do not reconfigure themselves without paying a significant coherence cost.

\subsection{Technical anchor: minimal stationary closure and correspondence principle}

For readers who wish to locate the foregoing interpretative claims within the formal PDL corpus, we briefly summarise the key technical ingredients.\cite{Laubscher_IntroPDL_2024,Laubscher_EmergencePhysicalReality_2024,Laubscher_MinimalClosures_46_2024,Laubscher_PDL_Position_2024}

First, the existence and necessity of the \((4,6)\) closure are established by a combinatorial analysis of finite signed graphs subject to the PDL axioms. One considers all graphs with a small number of vertices and imposes: (i) the possibility of a non-trivial logical \(2\)-cycle on edge signs; (ii) a high degree of triangular coherence according to a fixed sign-product closure condition; (iii) minimal completeness, in the sense that no implicit external structure is invoked to repair incoherences; and (iv) optimisation criteria that favour connectivity and coherence under strict combinatorial minimality.\cite{Laubscher_MinimalClosures_46_2024} The main result is that no graph with fewer than four vertices satisfies these constraints, and that among graphs with four vertices the complete graph with six edges admits sign assignments that render all triangles coherent while supporting a global sign-inversion \(2\)-cycle. This graph is unique up to isomorphism and sign relabelling under the stated constraints.\cite{Laubscher_MinimalClosures_46_2024}

Second, the correspondence with the electron is formulated not as a theorem but as an explicit hypothesis, articulated in the main technical paper on the emergence of physical reality within PDL.\cite{Laubscher_EmergencePhysicalReality_2024} The proposal is that the electron at rest realises, at the physical level, the logically defined \((4,6)\) closure, with its Compton frequency corresponding to the internal pulsation and the reduced Planck constant quantifying the coherence cost per cycle. In this setting, the relation \(E = \hbar \omega\) is reinterpreted as a statement about the energetic price of maintaining the minimal stationary regime defined combinatorially. This proposal is calibrated against known electron properties but remains open to refinement or refutation as the PDL framework and empirical constraints evolve.\cite{Laubscher_EmergencePhysicalReality_2024,Laubscher_PDL_Position_2024,Laubscher_PDL_GlobalMapping_v6_2024}

Third, the broader research programme document situates this correspondence within a hierarchy of further constructions: the combinatorial architecture of the proton, structural derivations of the fine-structure constant, treatments of coherence leakage and effective fields, and preliminary steps towards an Einstein--Dirac interface.\cite{Laubscher_PDL_Position_2024,Laubscher_ProtonArch_v2_2024,Laubscher_Alpha_v2_2024,Laubscher_EffectiveFields_2024,Laubscher_Einstein_Dirac_2024} The guiding methodological stance is conservative: whenever an interpretative step is taken, its status (established, conjectural, or speculative) is explicitly recorded. The present paper follows the same discipline, while allowing itself to dwell more extensively on the ontological and existential resonances of treating the \((4,6)\) closure as the first ``instrument'' of a logic of coherence.

\section{From Pulses to Architectures: Protons, Surfaces, and Constants}
\label{sec:architectures}

\subsection{Beyond the first closure: hierarchical architectures}

If the \((4,6)\) block represents the first admissible pulsating closure, the world we inhabit clearly requires more than copies of this minimal pattern. Atoms, molecules, macroscopic bodies, and fields all point to richer architectures in which many elementary closures are organised into composite structures. One of the central claims of PDL is that such organisation need not be imposed by hand: given the axioms and optimisation criteria already described, there are combinatorial reasons for why certain multi-scale configurations are singled out as especially coherent.\cite{Laubscher_EmergencePhysicalReality_2024,Laubscher_PDL_GlobalMapping_v6_2024}

The proton provides the first non-trivial test case for this claim. In the PDL framework, it is not modelled as a featureless point or as an unconstrained assembly of quarks, but as a hierarchically structured composite built from \((4,6)\) units arranged under additional axiomatic requirements on connectivity, symmetry, and closure quality.\cite{Laubscher_ProtonArch_v1_2024,Laubscher_ProtonArch_v2_2024} The outcome of this construction is surprisingly rigid. Instead of a continuum of possible configurations, one obtains a small set of integer-valued parameters---valence occupancies \(24\) and \(28\), a valence content \(r_{\mathrm{val}} = 930\), a relational sea with \(10\,087\) links, and a total relational budget of \(11\,017\)---that characterise a minimally complex stationary architecture whose internal coherence can be scrutinised and compared with empirical data.\cite{Laubscher_ProtonArch_v2_2024,Laubscher_PDL_GlobalMapping_v6_2024}

This hierarchical view of the proton reinforces the guiding sentence in a new register. If even a seemingly simple nucleon is the outcome of a delicate optimisation of coherence across several nested levels---core closures, relational seas, active interfaces---then the idea that ``whatever we may be, we are expressions of a relentless logic of coherence'' ceases to be an abstraction applied only to electrons. It becomes a concrete statement about how matter, in its first non-trivial bound form, is already an instrument through which that logic achieves a remarkably specific and constrained organisation.\cite{Laubscher_Existence_PulsatingClosure_2026}

\subsection{The golden ratio and the active surface}

Within this composite architecture, special attention is given in PDL to the notion of an \emph{active surface}: a finite set of relations through which the proton couples to external closures, notably to electron-like \((4,6)\) blocks.\cite{Laubscher_GoldenRatio_Sketch_2024,Laubscher_Born_Golden_v2_2024} The active surface is neither an arbitrary boundary nor a purely geometric shell. It is defined combinatorially as the subset of the proton’s relational budget that is effectively projected outward and made available for binding and interaction, while the remainder of the structure remains internally focused on maintaining its own coherence.

A striking feature of the PDL proton model is the appearance of the golden ratio \(\varphi\) in the optimisation of this surface. Under the imposed constraints, the partition between internal closure and external coupling is best realised, in a minimal sketch, by a core–surface decomposition whose proportions are governed by \(\varphi\).\cite{Laubscher_GoldenRatio_Sketch_2024,Laubscher_Alpha_v2_2024} This is not introduced as a symbolic flourish; it arises as a candidate solution to a concrete combinatorial problem: how to allocate a finite coherence budget between bulk and interface so as to maintain a stable architecture that can still interact non-trivially with its environment.

In this context, the role of the golden ratio remains conjectural but tightly constrained. It suggests that even the way a proton exposes itself to the rest of the world is not free, but shaped by optimisation criteria that balance the need for internal robustness against the need for communicative openness.\cite{Laubscher_GoldenRatio_Sketch_2024,Laubscher_Born_Golden_v2_2024} Philosophically, this lends additional substance to the idea that existence, in the PDL sense, is inseparable from having a regulated mode of exposure: to be is to maintain a boundary that neither collapses into isolation nor dissolves into uncontrolled leakage.

\subsection{Constants as coherence ratios}

One of the most distinctive aspects of PDL is its treatment of certain physical constants not as primitive numerical inputs, but as emergent ratios that quantify how coherence is distributed within and around composite closures.\cite{Laubscher_Alpha_v2_2024,Laubscher_Exponent18_v2_2024,Laubscher_UniversalCoherenceLeakage_2024} In the proton model, for example, the fine-structure constant \(\alpha\) is reinterpreted as a surface-to-volume coherence ratio: a measure of what fraction of the proton’s total coherence budget is effectively projected onto its active surface and thus available for electromagnetic coupling with an electron.\cite{Laubscher_Alpha_v2_2024,Laubscher_Born_Golden_v2_2024}

Similarly, Boltzmann’s constant can be associated with how coherence is partitioned between micro-level relational structures and macro-level thermodynamic descriptions, while an effective gravitational coupling emerges from the leakage of coherence into a surrounding relational network that extends beyond any single proton.\cite{Laubscher_UniversalCoherenceLeakage_2024,Laubscher_EffectiveFields_2024,Laubscher_PDL_Position_2024} In each case, the constant is not treated as a brute fact; it is regarded as a compact way of encoding the outcome of a multi-level coherence bookkeeping exercise.

From this viewpoint, constants become signatures of the way the universe manages its coherence budgets across scales. The fine-structure constant, for instance, can be seen as expressing a very specific compromise between internal stability and external responsiveness in the proton–electron system; the gravitational coupling reflects how much unresolved coherence must be exported into a global field-like structure when local closures have been pushed as far as possible.\cite{Laubscher_Alpha_v2_2024,Laubscher_UniversalCoherenceLeakage_2024} To the extent that these proposals withstand empirical and mathematical scrutiny, they reinforce the overarching picture: what we call the ``laws of nature'' may be read as concise descriptions of how a finite, logically constrained universe allocates coherence in order to remain both organised and open.

\subsection{Technical anchor: combinatorial proton and constant derivations}

Technically, the proton architecture in PDL is obtained by imposing a set of axiomatic constraints on assemblies of \((4,6)\) blocks: constraints on valence structure, connectivity, symmetry, and closure quality, together with optimisation criteria similar in spirit to those that single out the minimal electron prototype.\cite{Laubscher_ProtonArch_v1_2024,Laubscher_ProtonArch_v2_2024} These conditions drastically reduce the space of admissible configurations and select a small family of architectures characterised by the integer invariants mentioned above. The resulting structures are then analysed in terms of their internal coherence and interaction capacities.

The structural derivation of the fine-structure constant proceeds by relating the active surface of the proton to the total relational budget and by considering how an electron-like \((4,6)\) closure couples to this surface.\cite{Laubscher_Alpha_v2_2024,Laubscher_Born_Golden_v2_2024} The proposed expression for \(\alpha\) combines architectural integers with factors involving the golden ratio and mass ratios, and yields a numerical value close to the empirical fine-structure constant without arbitrary fitting. Related constructions link coherence leakage parameters and hierarchical filtering exponents to an effective gravitational coupling.\cite{Laubscher_Exponent18_v2_2024,Laubscher_UniversalCoherenceLeakage_2024}

In all such derivations, the PDL programme is explicit about the epistemic status of its claims. Some elements---such as the combinatorial rigidity of the proton architecture under the stated axioms---are established by detailed analysis.\cite{Laubscher_ProtonArch_v2_2024,Laubscher_PDL_GlobalMapping_v6_2024} Others---notably the precise expressions for constants---are presented as conjectural but tightly constrained proposals, open to refinement or falsification as both the mathematical framework and the empirical data evolve.\cite{Laubscher_Alpha_v2_2024,Laubscher_PDL_Position_2024} The present paper relies on this technical work only to the extent required to support its interpretative thesis: that physical constants can be read as compressed descriptions of how a finite universe distributes coherence between interior, boundary, and beyond.

\section{Incomplete Closure: Leakage, Transcendence, and Global Coherence}
\label{sec:leakage_transcendence}

\subsection{No perfect wholes: the necessity of leakage}

Thus far, closure and coherence have been treated as ideals: minimal pulsating structures that satisfy all their local constraints and composite architectures that optimise stability and interaction. Yet, one of the recurrent lessons of the PDL corpus is that such ideals are never fully realised at higher levels. Even when a composite structure has been tuned as carefully as possible, an irreducible fraction of its triangular closure conditions remains unsatisfied.\cite{Laubscher_Leakage_Probability_2024,Laubscher_Exponent18_v2_2024,Laubscher_PDL_Position_2024} This residual inconsistency is not a technical blemish to be corrected later; it is promoted to a central structural feature.

PDL quantifies this feature by introducing a leakage parameter \(\varepsilon\), defined as the ratio of violated triangles to the total number of triangles in the relevant constraint graph.\cite{Laubscher_Leakage_Probability_2024} For elementary closures such as the \((4,6)\) block, one can achieve exact coherence with \(\varepsilon = 0\).\cite{Laubscher_MinimalClosures_46_2024} For composite configurations---protons, effective fields, cavity modes---\(\varepsilon\) remains small but strictly positive, even after optimisation.\cite{Laubscher_ProtonArch_v2_2024,Laubscher_EffectiveFields_2024,Laubscher_CavityModes_PDL_2024} In other words, the more ambitious the closure, the more it must live with a controlled failure to be perfectly closed.

This observation has both technical and ontological weight. Technically, it underpins PDL accounts of probabilistic behaviour, effective field descriptions, and gravitational coupling.\cite{Laubscher_Leakage_Probability_2024,Laubscher_EffectiveFields_2024,Laubscher_UniversalCoherenceLeakage_2024} Ontologically, it challenges any image of the universe as composed of fully self-sufficient wholes. If complex structures cannot eliminate leakage without sacrificing their very complexity, then incompleteness is not an accident; it is the price of higher-order existence.

\subsection{From local defect to global coupling}

What becomes of the leaked coherence? PDL does not treat \(\varepsilon\) as a mere error term to be discarded, but as a signal that part of the coherence cost associated with a composite structure cannot be accommodated within that structure alone.\cite{Laubscher_Leakage_Probability_2024,Laubscher_UniversalCoherenceLeakage_2024} In the case of nucleonic architectures, for example, the leakage fraction can be interpreted as an intrinsic deficit that must be exported into a surrounding relational network. This exported coherence does not vanish; it reappears as a modification of the compatibility relations between closures and, at a coarse-grained level, as a contribution to an emergent metric or field.

In more technical terms, PDL introduces the idea of hierarchical filtering: local leakage at the proton scale, characterised by a small \(\varepsilon\), is propagated through a sequence of coherence filters, each associated with a higher organisational level.\cite{Laubscher_Exponent18_v2_2024,Laubscher_UniversalCoherenceLeakage_2024} Under an exponentiation process (often modelled with an effective exponent around \(18\)), an initially minute deficit is amplified into a macroscopic modulation of coherence cost that can be read as a gravitational coupling or as a contribution to effective field strengths.\cite{Laubscher_Exponent18_v2_2024,Laubscher_EffectiveFields_2024} What began as a local failure of closure thus becomes a global constraint on how other closures can coexist and move.

From this perspective, leakage is not merely tolerated; it is the very mechanism by which finite closures participate in a larger coherence field. A proton that tried to seal itself perfectly would, in the PDL picture, forgo the possibility of contributing to metric structure. Conversely, the existence of a smooth large-scale order is traced back to the fact that no composite entity succeeds in internalising all its coherence obligations.\cite{Laubscher_UniversalCoherenceLeakage_2024,Laubscher_PDL_Position_2024} The universe holds together, in part, because no part can fully close upon itself.

\subsection{Transcendence and the impossibility of total self-containment}

The notion of leakage lends itself naturally to a philosophical vocabulary that has often been reserved for theology or metaphysics. In collaboration and dialogue with the Theory of Objectivity, PDL has explored the idea that the impossibility of perfect closure can be read as a disciplined notion of transcendence.\cite{Laubscher_PDL_TO_Dialogue_2024,Laubscher_Existence_PulsatingClosure_2026} Here, transcendence does not mean a separate realm beyond the physical world; it denotes the structural fact that any finite closure necessarily refers beyond itself, because it cannot fully absorb the coherence requirements it generates.

In this sense, every composite structure exists in a mode of regulated exposure. It sustains a finite core of internal relations, maintains an active surface through which it interacts, and exports a fraction of its unresolved coherence into a wider network that, in turn, constrains it.\cite{Laubscher_ProtonArch_v2_2024,Laubscher_UniversalCoherenceLeakage_2024} No entity is entirely ``at home'' within its own boundaries; part of what it is depends on conditions that are instantiated elsewhere. The leakage parameter \(\varepsilon\) measures, in a starkly combinatorial way, the degree to which an entity’s existence is structurally incomplete and therefore open.

The guiding sentence of this paper can now be heard in a new register. To say that we are ``the most vulnerable instrument'' of a logic of coherence is to acknowledge that our own closures---biological, cognitive, social---are pervaded by leakage in exactly this sense.\cite{Laubscher_Existence_PulsatingClosure_2026} We cannot stabilise ourselves without relying on broader relational structures that we do not control, and yet we remain exposed to them. Vulnerability, on this reading, is not a contingent misfortune, but an expression of the same impossibility of total self-containment that already appears in the proton and in the emergent metric field. 

At the human scale, this structural openness is not an abstraction. Our bodies cannot maintain their organisation without a continuous exchange of matter and energy with an environment; our identities cannot stabilise without relying on shared languages, institutions, and memories that no single person controls. In PDL terms, our biological and social closures live with their own leakage parameters: there is always a part of what we are that is carried, sustained, or threatened by structures outside us. The mathematics of incomplete closure thus mirrors a very concrete aspect of lived existence—the impossibility of being entirely self-sufficient.

\subsection{Technical anchor: leakage, exponent 18, and emergent couplings}

Formally, leakage in PDL is introduced by defining a coherence functional on the space of signed graphs, typically counting satisfied versus violated triangular closure conditions.\cite{Laubscher_Leakage_Probability_2024} For a given architecture, one computes \(\varepsilon\) as the fraction of violated triangles at the optimal sign assignment. This parameter enters several subsequent constructions. In probabilistic reformulations, leakage supports interpretations of probabilities as self-maintained frequencies of closure violations.\cite{Laubscher_Leakage_Probability_2024} In the context of effective fields, it governs how local deficits of coherence are re-expressed as contributions to field-like degrees of freedom.\cite{Laubscher_EffectiveFields_2024}

The role of the exponent \(18\) arises in models where leakage is filtered through a hierarchy of organisational levels.\cite{Laubscher_Exponent18_v2_2024} By iterating a modest amplification factor across multiple coherence filters, one can map a tiny nucleonic \(\varepsilon\) onto a macroscopic coupling strength compatible with observed gravitational interactions.\cite{Laubscher_UniversalCoherenceLeakage_2024} These constructions are explicitly presented as conjectural, with clear criteria for their refinement or rejection.\cite{Laubscher_PDL_Position_2024,Laubscher_PDL_GlobalMapping_v6_2024}

For the purposes of the present paper, what matters is less the specific numerical value of the exponent than the structural pattern it encodes: leakage is never simply thrown away. It is systematically transformed into a constraint that binds together otherwise distinct closures. In this sense, the technical machinery of \(\varepsilon\) and its hierarchical filtering provides a concrete backbone for the philosophical claim that no finite entity is ontologically self-sufficient, and that what might be called ``transcendence'' is inscribed in the very grammar of closure.

\section{Causality Reframed: From Chains to Coherent Patterns}
\label{sec:causality}

\subsection{Classical chains and their limits}

In classical mechanics, causality is often pictured as a chain of events linked by deterministic laws: given the state of a system at one time, the laws dictate its state at all later times. This image, sometimes reinforced by metaphors of clockwork universes, suggests that causality is nothing more than the successive application of local rules on a pre-given spacetime stage. While this picture has been remarkably successful in many regimes, it becomes conceptually fragile at the extremes that motivated PDL in the first place: near singularities, in quantum field theory, or in hypothetical cosmological origins.\cite{Laubscher_EmergencePhysicalReality_2024,Laubscher_PDL_Position_2024}

In such regimes, it is not just that we lack data or computational power; the very notions of state, time, and trajectory become ambiguous. When spacetime can no longer be modelled as a smooth manifold, or when particle identities lose their usual clarity, it is legitimate to ask whether the classical image of causality as a linearly ordered chain of events is still appropriate.\cite{Laubscher_Existence_PulsatingClosure_2026} PDL responds to this question by stepping back from the continuum and rephrasing causality in terms of global coherence on finite relational structures.

\subsection{Global coherence as causal order}

In the PDL framework, the primary object is not a trajectory in spacetime, but a pulsating closure on a signed graph.\cite{Laubscher_IntroPDL_2024,Laubscher_EmergencePhysicalReality_2024} The dynamics of interest is the repeated application of an update rule that reassigns edge signs while striving to maintain or optimise triangular coherence. From this vantage point, what matters is not the local succession of configurations as such, but whether the pattern of pulsation remains compatible with the closure constraints that define the structure.

Causality can then be reinterpreted as the way in which a global requirement of coherence is realised across a network of relations, rather than as a mere sequence of pushes and pulls. An ``event'' in this setting is a change in the sign configuration that either preserves the viability of the closure or drives it towards breakdown.\cite{Laubscher_Leakage_Probability_2024,Laubscher_EffectiveFields_2024} To say that one event causes another is to say that the second is required if the global coherence budget is to be managed under the constraints fixed by the first. The order here is not primarily temporal; it is structural: certain configurations can coexist at low coherence cost, others cannot.\cite{Laubscher_Existence_PulsatingClosure_2026}

This perspective does not deny local regularities. It rather embeds them in a broader logic: what appears, in a coarse-grained description, as a law relating local states may emerge from the deeper requirement that certain closure patterns be globally maintainable on finite graphs.\cite{Laubscher_PDL_GlobalMapping_v6_2024} The guiding sentence’s reference to a ``logic of causality that fulfils itself with every pulse'' is thus given a precise meaning: each pulsation is an occasion on which the system either succeeds or fails in reasserting a globally coherent relational order.

\subsection{Constants as structural signatures of causality}

Within this reframed view, physical constants acquire a natural causal interpretation. If constants such as the fine-structure constant, Boltzmann’s constant, or effective gravitational couplings encode how coherence is apportioned between interior, surface, and leakage, then they also encode which evolutions are structurally admissible.\cite{Laubscher_Alpha_v2_2024,Laubscher_UniversalCoherenceLeakage_2024,Laubscher_PDL_Position_2024} A given value of \(\alpha\), for example, does not merely describe the strength of electromagnetic interactions; it specifies how tightly an electron-like closure can bind to a proton surface without destabilising the multi-level architecture on which both depend.\cite{Laubscher_Alpha_v2_2024,Laubscher_Born_Golden_v2_2024}

In this sense, constants become signatures of a deeper causal order: they summarise, in a few numbers, the conditions under which closures can coexist and interact without catastrophic coherence loss. The ``laws of nature'' are then read as concise statements about the management of coherence budgets across scales, rather than as arbitrary rules imposed on an inert substrate.\cite{Laubscher_PDL_GlobalMapping_v6_2024} To say that the universe obeys these laws is to say that its pulsating closures repeatedly succeed in satisfying the associated coherence constraints.

\subsection{Technical anchor: coherence cost, mode structures, and emergent metrics}

Technically, several strands of PDL work illustrate how global coherence requirements give rise to structures that resemble familiar causal orderings. The reinterpretation of \(E = \hbar \omega\) as a statement about the coherence cost per cycle ties energy and frequency directly to the maintenance of minimal closures.\cite{Laubscher_EmergencePhysicalReality_2024,Laubscher_Schrodinger_Compat_v2_2024} The analysis of discrete cavity modes demonstrates that even in the absence of a continuum background, signed graphs can support mode families indexed by integer triples, with periods and densities of states that mirror those of standard field-theoretic cavities.\cite{Laubscher_CavityModes_PDL_2024}

Similarly, work on effective fields and coherence leakage shows how local deficits can be re-expressed as contributions to emergent metric structures and couplings.\cite{Laubscher_EffectiveFields_2024,Laubscher_UniversalCoherenceLeakage_2024} Taken together, these constructions suggest that what we call causal structure in spacetime may be, at least in part, a coarse-grained reflection of how coherence constraints organise pulsating closures on finite graphs. The present paper does not attempt to complete this reconstruction, but uses these results as anchors for its claim that causality can be understood as global coherence in action.

\section{Freedom, Vulnerability, and Human Existence}
\label{sec:freedom_human}

\subsection{The emergence of freedom from necessity}

If the universe is governed by a strict logic of coherence, it may seem that there is little room left for freedom. A naïve reading of PDL might suggest a form of structural determinism: once the axioms and optimisation criteria are fixed, everything else follows. Yet the PDL corpus itself counsels against such a conclusion. It repeatedly distinguishes between what is strictly derived, what is conjectural, and what remains open, and it explicitly warns against reducing complex phenomena to their microphysical underpinnings.\cite{Laubscher_Existence_PulsatingClosure_2026,Laubscher_PDL_Position_2024}

Within this cautious framework, freedom can be approached not as a violation of physical regularities, but as a high-level expression of the same logic of closure and leakage that governs protons and fields. Complex biological and cognitive systems can be viewed, in PDL-friendly terms, as high-order closures that have internalised part of their own coherence management: they allocate, monitor, and reallocate coherence budgets across many interacting substructures.\cite{Laubscher_Existence_PulsatingClosure_2026} When such a system becomes capable of representing aspects of its own organisation and of exploring alternative ways of maintaining itself, its behaviour may legitimately be described as involving choice and deliberation, even if it never steps outside the logic of coherence.

In this view, necessity and freedom are not opposites but successive expressions of the same underlying project. The same constraints that fix the architecture of the proton also make possible the emergence of nervous systems and brains. Once these higher-level closures arise, they can inherit the logic that produced them and bend it back upon itself, exploring variations and novel arrangements that were not rigidly prescribed by micro-level structures.\cite{Laubscher_PDL_TO_Dialogue_2024} Freedom, then, is not the absence of structure; it is what a deeply structured universe can do when its own closures acquire the capacity to reorganise themselves.

\subsection{Humans as vulnerable instruments}

The phrase ``the human being is nothing more than its most vulnerable instrument'' can now be unpacked in light of leakage and incomplete closure. Human organisms and minds are not isolated monads; they are open, multi-layered closures whose persistence depends on a delicate interplay between internal stability and external support.\cite{Laubscher_Existence_PulsatingClosure_2026} They must continuously negotiate boundaries: which exchanges of matter, energy, and information they can afford; which they must resist; which they must seek out in order to maintain or expand their coherence.

Vulnerability, in this context, is not merely biological fragility or psychological sensitivity. It is the structural condition that follows from incomplete closure. Just as a proton cannot exist as a fully sealed object but must export part of its coherence into a broader relational field, so too a person cannot exist without relying on infrastructures, ecosystems, social relations, and symbolic networks that extend far beyond any individual control.\cite{Laubscher_UniversalCoherenceLeakage_2024,Laubscher_PDL_TO_Dialogue_2024} To be an ``instrument'' of the logic of coherence is to be a site where that logic is both realised and tested, under conditions of partial exposure and inevitable dependence.

From this angle, many familiar aspects of human life can be re-read in PDL terms. The effort required to maintain an identity over time, the strain of conflicting commitments, the relief brought by resolving long-standing contradictions, all resemble the dynamics of closures that struggle to keep their coherence budget within sustainable bounds.\cite{Laubscher_Existence_PulsatingClosure_2026} Vulnerability is the openness that makes such struggles meaningful: without it, there would be no possibility of growth, learning, or ethical responsibility, but also no genuine risk.

\subsection{Coherence budgets in lived experience}

Although PDL is not a theory of psychology, its vocabulary of coherence, closure, leakage, and active surfaces offers suggestive analogies for thinking about lived experience.\cite{Laubscher_Existence_PulsatingClosure_2026} One might, with care, compare attention to the allocation of a limited coherence budget across competing tasks; identity to the maintenance of a relatively stable pattern of closures across changing contexts; memory to the retention of structural regularities long enough for them to be re-instantiated and shared.

These analogies are not intended to reduce mental life to graph theory. They are invitations to notice that many of our most intimate difficulties and achievements can be described in terms strikingly similar to those used for complex closures. When someone tries to reconcile incompatible commitments, or to integrate a painful experience without losing their sense of self, they are effectively attempting to extend their pattern of coherence while accepting a certain amount of controlled leakage. When they fail—when contradictions multiply beyond what their current structure can accommodate—the result is experienced as confusion, breakdown, or crisis. In this light, PDL does not tell us what to feel or how to live, but it offers a structural lens for understanding why the work of maintaining a life that “hangs together” can be both fragile and demanding.

\subsection{Methodological caution and non-reduction}

Because these correspondences are tempting, it is essential to restate the methodological caution that runs through the PDL research programme.\cite{Laubscher_PDL_Position_2024,Laubscher_Existence_PulsatingClosure_2026} PDL was developed primarily as a proposal about the organisation of microphysical reality and the emergence of constants and effective fields. When it is extended towards life, mind, or meaning, this extension is explicitly labelled as speculative and analogical. No derivation of phenomenology is claimed from first principles on signed graphs.

The ambition of the present paper is therefore modest in one respect and ambitious in another. Modest, because it refuses to overstate what PDL can currently say about human existence; ambitious, because it insists that even carefully limited analogies can change the way we understand our place in the universe. If reality is indeed organised as a hierarchy of pulsating, incomplete closures that negotiate coherence and leakage, then our own vulnerability and capacity for freedom are not afterthoughts tacked onto an indifferent world. They are particular ways in which the underlying logic of coherence becomes, in the guiding sentence’s words, an engine and an essence that finds in us one of its most exposed instruments.

\section{Positioning PDL: Realism, Structuralism, and Information}
\label{sec:positioning}

\subsection{Relation to ontic structural realism}

From a philosophical point of view, PDL can naturally be read as a concrete candidate for ontic structural realism (OSR). OSR holds, in broad terms, that what fundamentally exists are not individual objects with intrinsic natures, but relational structures, patterns, or networks of relations.\cite{Laubscher_Existence_PulsatingClosure_2026,Laubscher_PDL_Position_2024} PDL embraces this orientation explicitly: its primitive entities are finite signed graphs, its basic bearers of existence are pulsating closures on those graphs, and its account of particles and fields is formulated in structural, rather than object-based, terms.\cite{Laubscher_EmergencePhysicalReality_2024,Laubscher_MinimalClosures_46_2024}

There are, however, important differences between PDL and many abstract discussions of structural realism. First, PDL is relentlessly finite. It eschews idealisations involving infinite lattices or continuous fields at the fundamental level, insisting that genuine existence must be realisable as a self-sustaining closure on a bounded configuration.\cite{Laubscher_EmergencePhysicalReality_2024,Laubscher_MinimalClosures_46_2024} Second, it does not rest content with the slogan that ``only structure exists.'' It provides explicit combinatorial mechanisms by which minimal closures, composite architectures, constants, and effective fields arise from its primitive structures.\cite{Laubscher_PDL_GlobalMapping_v6_2024} In this sense, PDL can be seen as an attempt to operationalise OSR: to turn a general metaphysical thesis into a detailed research programme.

\subsection{Relation to informational and digital ontologies}

At the same time, PDL invites comparison with informational or digital ontologies, in which the fundamental level of reality is described in terms of bits, qubits, or abstract computational processes. It shares with such approaches the intuition that relational and combinatorial organisation precedes familiar physical quantities, and that much of physics can be reconstructed from constraints on information-like structures.\cite{Laubscher_EmergencePhysicalReality_2024,Laubscher_EffectiveFields_2024} Yet it resists two common tendencies in these narratives.

First, PDL is resolutely ontic rather than epistemic. Its signed graphs are not conceived as records of measurements or as representations in an observer’s mind; they are proposed as the very fabric of what exists.\cite{Laubscher_PDL_Position_2024,Laubscher_Leakage_Probability_2024} Coherence is not a norm governing rational belief, but a structural condition for the persistence of closures. Second, PDL insists on mathematical discipline. Where some informational metaphors remain vague about what counts as a state, an update, or a closure, PDL offers precise definitions, axioms, and theorems.\cite{Laubscher_IntroPDL_2024,Laubscher_MinimalClosures_46_2024} In this respect, it is less a poetic ``it from bit'' story than a structured attempt to say exactly how much and what kind of relational information is needed for existence.

\subsection{Interfaces with other foundational approaches}

Beyond OSR and informational views, PDL also interacts with other foundational programmes, notably the Theory of Objectivity (TO).\cite{Laubscher_PDL_TO_Dialogue_2024,Laubscher_PDL_Position_2024} TO emphasises the role of boundaries, multi-observer stabilisation, and a transcendent substance that exceeds any particular quantum state. PDL provides a complementary technical apparatus: boundaries are instantiated as active surfaces and finite interfaces; multi-observer stabilisation finds an analogue in the coexistence and mutual compatibility of closures; and transcendence is modelled via unavoidable leakage into a broader relational field.\cite{Laubscher_Leakage_Probability_2024,Laubscher_UniversalCoherenceLeakage_2024}

These correspondences do not imply identity between the frameworks. Rather, they suggest that PDL can serve as a concrete playground in which concepts drawn from TO and related approaches can be tested, sharpened, or revised.\cite{Laubscher_PDL_TO_Dialogue_2024} More generally, PDL’s emphasis on finitely realised closures and explicit coherence measures makes it a natural interlocutor for any foundational theory that takes seriously the interplay between local structure, global constraints, and the limits of self-containment.

\subsection{Technical anchor: global mapping and research programme}

The overall shape of the PDL research programme, and the status of its various components, have been charted in a series of global mapping and position papers.\cite{Laubscher_PDL_Position_2024,Laubscher_PDL_GlobalMapping_v6_2024} These documents classify results into categories: rigorously established constructions (such as the existence and necessity of the \((4,6)\) block), strongly constrained architectural proposals (such as the proton structure and associated constants), and speculative extensions (such as effective fields, gravitational couplings, and ontological reflections on existence and transcendence).\cite{Laubscher_MinimalClosures_46_2024,Laubscher_ProtonArch_v2_2024,Laubscher_Alpha_v2_2024,Laubscher_UniversalCoherenceLeakage_2024,Laubscher_Existence_PulsatingClosure_2026}

This explicit epistemic bookkeeping is part of what justifies treating PDL as a research programme rather than as a dogmatic system. It offers not only hypotheses about the structure of reality, but also clear indications of how these hypotheses might be refined, constrained, or falsified.\cite{Laubscher_PDL_Position_2024,Laubscher_PDL_GlobalMapping_v6_2024} The present paper fits within this cartography as a synthetic and interpretative piece: it does not expand the technical map, but attempts to show how the existing landmarks can be read together as a coherent ontological proposal.

\section{A Logically Constrained Universe: Risks, Tests, and Horizons}
\label{sec:risks_tests}

\subsection{What must existence be like, if PDL is right?}

The question that animates PDL can be phrased succinctly: \emph{what must existence be like if a universe qualitatively like ours is to arise at all?} Rather than starting from observed regularities and seeking ever more refined differential equations, PDL asks what minimal logical and structural conditions must be met for there to be persistent structures, constants, and effective fields in the first place.\cite{Laubscher_EmergencePhysicalReality_2024,Laubscher_IntroPDL_2024,Laubscher_PDL_Position_2024}

The answer it proposes is both simple and demanding. Existence, on this view, consists in being a pulsating closure on a finite signed graph, selected and sustained by optimisation of coherence under constraint. Electrons, protons, fields, and metrics are not primitive givens but higher-order expressions of this logic: ways in which minimal closures compose into complex architectures, allocate coherence budgets, and negotiate inevitable leakage into broader relational networks.\cite{Laubscher_MinimalClosures_46_2024,Laubscher_ProtonArch_v2_2024,Laubscher_EffectiveFields_2024,Laubscher_UniversalCoherenceLeakage_2024} If this picture is even approximately right, then the guiding sentence of this paper is more than a metaphor; it is an attempt to articulate in everyday language what the formalism already says.

\subsection{Potential objections and failure modes}

A programme as ambitious as PDL naturally invites objections. One cluster of concerns targets underdetermination: could other discrete or relational frameworks reproduce the same numerical successes without sharing PDL’s ontological commitments? Another cluster questions the reliance on specific architectural integers, leakage exponents, or optimisation criteria: are these truly forced by minimality, or might they conceal hidden tuning?\cite{Laubscher_PDL_Position_2024,Laubscher_PDL_GlobalMapping_v6_2024}

There are also worries about scope. Even if PDL offers a compelling account of certain constants or microstructures, does it have the resources to address the full range of phenomena encompassed by contemporary physics, from cosmology to condensed matter? And, as this paper itself illustrates, there is the ever-present risk of overextending structural metaphors into domains where they are suggestive but not yet well-grounded.\cite{Laubscher_Existence_PulsatingClosure_2026} Any honest assessment of PDL must take these potential failure modes seriously, not as reasons to dismiss the programme out of hand, but as guides for where it must be strengthened, constrained, or perhaps abandoned.

\subsection{Empirical and mathematical tests}

PDL is not immune to empirical and mathematical scrutiny. On the empirical side, its proposals linking proton architecture to constants such as \(\alpha\), or linking leakage to gravitational coupling, can be refined and tested against high-precision measurements.\cite{Laubscher_Alpha_v2_2024,Laubscher_UniversalCoherenceLeakage_2024} Discrepancies would not automatically falsify the entire framework, but they would force revisions of specific architectural assumptions or optimisation principles. On the mathematical side, the internal consistency and uniqueness claims of the architecture constructions can be challenged by constructing alternative graph models that satisfy the same axioms but yield different numerical predictions.\cite{Laubscher_ProtonArch_v2_2024,Laubscher_MinimalClosures_46_2024}

More broadly, the PDL programme can be evaluated by its fertility. Does it continue to generate non-trivial connections between discrete coherence structures and known physical regularities? Does it offer new questions, concepts, or tools that remain useful even if some of its specific conjectures fail?\cite{Laubscher_PDL_GlobalMapping_v6_2024,Laubscher_CavityModes_PDL_2024} The present paper does not attempt to settle these assessments, but assumes that PDL’s ultimate value will depend less on any one derivation than on the overall coherence and testability of the picture it offers.

\subsection{Open horizons: life, mind, meaning}

Beyond physics, PDL opens horizons rather than closing debates. Its account of existence as pulsating closure, of boundaries as active surfaces, and of transcendence as unavoidable leakage suggests a way of thinking about life, mind, and meaning that neither ignores physical structure nor collapses it into reduction.\cite{Laubscher_Existence_PulsatingClosure_2026,Laubscher_PDL_TO_Dialogue_2024} Living systems can be tentatively pictured as closures that have learned to manage their coherence budgets in their own favour; conscious experience as a regime in which certain closures can represent and evaluate aspects of their internal organisation; meaning as the emergence of patterns that stabilise across many layers of closure and exposure.

These suggestions are intentionally cautious. They do not claim to derive biology, psychology, or ethics from PDL, but to hint at how an ontology of finite, pulsating, incomplete closures might resonate with these domains.\cite{Laubscher_Existence_PulsatingClosure_2026} Whether this resonance proves fruitful will depend on future work at the interface of physics, philosophy, and the human sciences. For now, it is enough to note that a logically constrained universe, as sketched here, is not a cold and indifferent backdrop. It is a place where the effort to maintain coherence under constraint is written into the fabric of reality, from electrons to thoughts.

\section{Conclusion: Learning to Read the Logos}
\label{sec:conclusion}

\subsection{From technical corpus to shared language}

The PDL framework was not conceived primarily as a philosophical doctrine; it grew out of a technical attempt to reformulate microphysics on finite signed graphs, derive minimal closures, and reinterpret constants as coherence ratios.\cite{Laubscher_EmergencePhysicalReality_2024,Laubscher_MinimalClosures_46_2024,Laubscher_Alpha_v2_2024,Laubscher_UniversalCoherenceLeakage_2024} Yet, as the corpus expanded, it became increasingly clear that these constructions carry an implicit ontological picture: reality as a hierarchy of pulsating closures, incomplete yet coherent, negotiating their own persistence through delicate allocations of coherence budget.

This paper has tried to give that picture a shared language. It has traced a path from the minimal \((4,6)\) closure to the proton and its active surface, from constants as coherence ratios to leakage and emergent metrics, from global coherence as causal order to freedom and vulnerability in human existence.\cite{Laubscher_ProtonArch_v2_2024,Laubscher_EffectiveFields_2024,Laubscher_CavityModes_PDL_2024,Laubscher_Existence_PulsatingClosure_2026} The guiding sentence that opened the text was not intended as an independent metaphysical claim, but as an attempt to say plainly what a technically demanding framework already implies about what it means to exist.

\subsection{An invitation rather than a doctrine}

Throughout, care has been taken to respect the internal epistemic discipline of the PDL programme: to distinguish between established results, constrained conjectures, and speculative extrapolations.\cite{Laubscher_PDL_Position_2024,Laubscher_PDL_GlobalMapping_v6_2024} PDL is presented not as a finished doctrine that demands assent, but as a research programme and an invitation. Its axioms, constructions, and interpretations are offered as hypotheses to be tested, criticised, refined, or replaced.

If the programme proves resilient, it may eventually justify reading the universe as a logically constrained project in which coherence, closure, and leakage play the roles once reserved for substance and law. If it fails, it may still leave behind a set of concepts and questions that sharpen our understanding of what is at stake when we speak of existence, causality, or transcendence. In either case, the value of PDL lies less in the comfort of its answers than in the clarity of the problems it dares to pose.

\subsection{Whoever we may be, again}

We can now return, finally, to the sentence that served as our point of departure. To say that ``whoever we may be, whatever we may be, we are the necessary expression of a simple and relentless logic of coherence'' is, in PDL’s terms, to acknowledge that our being depends on the success of innumerable closures---from \((4,6)\) blocks and proton architectures to neural assemblies and social structures---in sustaining themselves under finite constraints.\cite{Laubscher_EmergencePhysicalReality_2024,Laubscher_ProtonArch_v2_2024,Laubscher_Existence_PulsatingClosure_2026} To call this logic an ``Origin, engine, and essence'' is to suggest that without it there would be nothing to speak of; with it, there can be a universe capable of generating vulnerable instruments such as us. 

On a daily scale, this may be as simple as recognising that every time we hold a thought steady, keep a promise, or repair a damaged relationship, we are participating in the same project of maintaining coherence under constraint that PDL attributes to electrons and protons. The scale and language change, but the underlying task—to preserve a pattern without freezing it, to stay open without dissolving—is structurally the same.

Whether one finds this picture comforting or unsettling may vary. What PDL offers is not consolation, but a framework within which such reactions can themselves be seen as part of the story: as responses of finite closures that have, for a brief time, become capable of reflecting on the conditions of their own existence. To learn to read the Logos, in this sense, is not to escape our vulnerability, but to understand it as one of the ways in which a logic of coherence looks back at itself through the eyes, and the questions, of whoever we may be.


\bibliographystyle{plain}
\bibliography{PDL_master_bib}

\end{document}
